%!TEX TS-program = lualatex
%!TEX encoding = UTF-8 Unicode

\documentclass[12pt, addpoints]{exam}

%\printanswers

\usepackage[left=1in,right=0.75in,top=1in]{geometry}     

\usepackage{fontspec}
\def\mainfont{Linux Libertine O}
\setmainfont[Ligatures={TeX, Common}, BoldFont={* Bold}, ItalicFont={* Italic}, Numbers={Proportional, OldStyle}]{\mainfont}
%\setmonofont[Scale=MatchLowercase]{Inconsolata} 
\setsansfont[Scale=MatchLowercase]{Linux Biolinum O} 
\usepackage{microtype}

% This defines \amper for the fancy ampersand
% to be used in the header. See
% https://tex.stackexchange.com/a/58185/39194
\usepackage{xspace}
\newfontfamily\amperfont[Style=Alternate]{Linux Libertine O}
\makeatletter
\DeclareRobustCommand{\amper}{{\amperfont\ifx\f@shape\scname\smaller[1.2]\fi\&}\xspace}
\makeatother

\usepackage{graphicx}
	\graphicspath{%
	{/Users/goby/Pictures/teach/466/exams/}}%

% Used to get the last two digits of the year for the header below.
\def\short#1{\csname @gobbletwo\expandafter\endcsname\number#1}

\pagestyle{headandfoot}
\firstpageheader{Ornithology, Exam 2, \textsc{s}\short{\year}, \numpoints~points.}{}{%
	\ifprintanswers\textbf{KEY}\else Name: \enspace \makebox[2.5in]{\hrulefill}\fi}
\runningheader{}{}{\small{pg. \thepage}}
\runningheadrule

\footer{}{}{\iflastpage{\small \rule{0.6in}{0.4pt} / \numpoints\ points}{}}%{\makebox[0.75in]{\hrulefill}}
\extrafootheight{-0.5in}

\pointsinmargin
\marginpointname{ pts}

\usepackage{multicol}

%% Remove comment %% to print answer key.
\unframedsolutions
\renewcommand{\solutiontitle}{}
\SolutionEmphasis{\bfseries}

%% Create a Matching question format
\newcommand*\Matching[1]{
\ifprintanswers
	\textbf{#1}
\else
	\rule{2in}{0.5pt}
\fi
}
\newlength\matchlena
\newlength\matchlenb
\settowidth\matchlena{\rule{2.1in}{0pt}}
\newcommand\MatchQuestion[2]{%
	\setlength\matchlenb{0.98\linewidth}
	\addtolength\matchlenb{-\matchlena}
	\parbox[t]{\matchlena}{\Matching{#1}}\enspace\parbox[t]{\matchlenb}{#2}}


\newcommand*\AnswerBox[2]{%
    \parbox[t][#1]{0.92\textwidth}{%
    \begin{solution}#2\end{solution}}
    \vspace{\stretch{1}}
}

\newenvironment{AnswerPage}[1]
    {\begin{minipage}[t][#1]{0.92\textwidth}%
    \begin{solution}}
    {\end{solution}\end{minipage}
    \vspace{\stretch{1}}}

\newlength{\basespace}
\setlength{\basespace}{5\baselineskip}

\begin{document}

\begin{questions}

\fullwidth{The first three questions require an audio recording. I'll play the recording for everyone. Once you have written your answers, you may complete the rest of the exam.}

\question
Answer the first three questions based on an audio recording of an Eastern Towhee. All I ask is that you try to do your best.

\begin{parts}
\part[4]
How many song types are heard in this recording? Enter the number below.

\AnswerBox{2\baselineskip}{3}

\part[3]
How many phrases are heard in the second song? Enter the number below.

\AnswerBox{2\baselineskip}{3 or 4}

\part[3]
How many notes are in the second phrase of the second song?

\AnswerBox{2\baselineskip}{1}

\end{parts}


\question[5]
Briefly describe the difference between a lek and an exploded lek.

\AnswerBox{0.2\textheight}{Answer}

\question[5]
What does an honest indicator tell the female about potential mates.

\AnswerBox{0.2\textheight}{Answer}

\newpage

\question[5]
Briefly explain why resource-based polygyny might evolve in birds?

\AnswerBox{0.2\textheight}{%
Answer goes here.}%
	
\question[10]
Eriefly explain the difference between the good genes hypothesis and the sexy sons hypothesis.

\AnswerBox{0.4\textheight}{Answer goes here.}

\newpage

\question[10]
Describe the functional difference between song and call in birds. 

\AnswerBox{0.2\textheight}{Answer goes here.}

\question[10]
Name the structure used by oscine birds to produce song. Describe how the structure works.

	\AnswerBox{0.2\textheight}{%
	Answer goes here.}

\newpage

%\question[10]
%Explain how having producer (finder) and scrounger individuals leads to an evolutionary stable strategy that maintains both foraging behaviors in the flock.
%	\AnswerBox{0.6\textheight}{Answer goes here.}%

\question[10]
Describe the difference between a sensitive period and opened ended learning for bird songs.

\AnswerBox{0.2\textheight}{Answer goes here.}%


\question[10]
Compare and contrast resource-based and mating territories. What do they have in common? What is different? 

\AnswerBox{0.2\textheight}{Answer goes here.}

	


\end{questions}


\end{document}